%% Copyright 2006-2015 Xavier Danaux (xdanaux@gmail.com).
%
% This work may be distributed and/or modified under the
% conditions of the LaTeX Project Public License version 1.3c,
% available at http://www.latex-project.org/lppl/.
%\documentclass[11pt,a4paper,sans]{moderncv}        % possible options include font size ('10pt', '11pt' and '12pt'), paper size ('a4paper', 'letterpaper', 'a5paper', 'legalpaper', 'executivepaper' and 'landscape') and font family ('sans' and 'roman')
%\usepackage{mathpazo}


\documentclass[11pt,a4paper]{moderncv}
\usepackage{mathpazo}

\usepackage[sfdefault,scaled=.9]{FiraSans}
\usepackage[T1]{fontenc}
\renewcommand*\oldstylenums[1]{{\firaoldstyle{} #1}}


\usepackage[utf8]{inputenc}
\usepackage{eurosym}

% moderncv themes
\moderncvstyle{classic}                             % style options are 'casual' (default), 'classic', 'banking', 'oldstyle' and 'fancy'
\moderncvcolor{blue}                               % color options 'black', 'blue' (default), 'burgundy', 'green', 'grey', 'orange', 'purple' and 'red'
%\renewcommand{\familydefault}{\sfdefault}         % to set the default font; use '\sfdefault' for the default sans serif font, '\rmdefault' for the default roman one, or any tex font name
%\nopagenumbers{}                                  % uncomment to suppress automatic page numbering for CVs longer than one page

\moderncvtheme{classic}
\usepackage[utf8]{inputenc}
% adjust the page margins
\usepackage[scale=0.8]{geometry}
\AtBeginDocument{\recomputelengths}

% personal data
\firstname{Bastian}
\familyname{Hagedorn}
\address{Einsteinstraße 62}{48149 Münster, Germany}
%\phone[fixed]{+49~(0)~251~83-32746}
\email{hagedornbastian@gmail.com}

\usepackage[style=numeric-verb,
            sorting=none, % keep order as in the bib file ...
            giveninits=true,
            defernumbers,
            backend=biber,
maxbibnames=50]{biblatex}

\defbibenvironment{bibliography}
  {\list%
     {\printfield{year}\hspace{1em}\printtext[labelnumberwidth]{\printfield{labelprefix}\printfield{labelnumber}}}
     {\setlength{\topsep}{0pt}% layout parameters from moderncvstyleclassic.sty
      \setlength{\labelwidth}{\hintscolumnwidth}%
      \setlength{\labelsep}{\separatorcolumnwidth}%
      \leftmargin\labelwidth%
      \advance\leftmargin\labelsep%
      }%
      \sloppy\clubpenalty4000\widowpenalty4000}
  {\endlist}
  {\item}


\DeclareNameAlias{default}{first-last}

\usepackage{xstring}
\usepackage{xpatch}
\newbibmacro*{name:bold}[2]{%
  \def\do##1{\iffieldequalstr{hash}{##1}{\bfseries\listbreak}{}}%
  \dolistloop{\boldnames}%
}

\newcommand*{\boldnames}{}

\xpretobibmacro{name:family}{\begingroup\usebibmacro{name:bold}{#1}{#2}}{}{}
\xpretobibmacro{name:given-family}{\begingroup\usebibmacro{name:bold}{#1}{#2}}{}{}
\xpretobibmacro{name:family-given}{\begingroup\usebibmacro{name:bold}{#1}{#2}}{}{}
\xpretobibmacro{name:delim}{\begingroup\normalfont}{}{}

\xapptobibmacro{name:family}{\endgroup}{}{}
\xapptobibmacro{name:given-family}{\endgroup}{}{}
\xapptobibmacro{name:family-given}{\endgroup}{}{}
\xapptobibmacro{name:delim}{\endgroup}{}{}

% Got hashes from the bbl file
\renewcommand*{\boldnames}{}
\forcsvlist{\listadd\boldnames}
  {{f3b57123791f63857d478ac84c802258},
   {7dd9b93410e4a44439b25ea5e2c49d4a}}

% Only print a year once
\newcounter{currentYear}
\DeclareFieldFormat{year}{%
\ifthenelse{\equal{#1}{\arabic{currentYear}}}%
    {}
%{\setcounter{currentYear}{#1}{\bfseries #1}}}
{\setcounter{currentYear}{#1}{#1}}}

\bibliography{publications}

\definecolor{color1}{HTML}{76B900}% NVIDIA Green
\definecolor{color2}{HTML}{9D9D9D}% NVIDIA Dark Grey
\definecolor{color3}{HTML}{008471}% NVIDIA Dark Green
\definecolor{color4}{HTML}{0071C5}% NVIDIA Blue

\newcommand{\strong}[1]{\textcolor{color0}{\textbf{#1}}}
\newcommand{\highlight}[1]{\textcolor{color0}{#1}}

\makeatletter
\renewcommand*{\cventry}[7][.25em]{%
  \cvitem[#1]{#2}{%
    {\strong{#3}}%
    \ifthenelse{\equal{#4}{}}{}{, {\slshape#4}}%
    \ifthenelse{\equal{#5}{}}{}{, #5}%
    \ifthenelse{\equal{#6}{}}{}{, #6}%
    .\strut%
    \ifx&#7&%
    \else{\newline{}\begin{minipage}[t]{\linewidth}\small#7\end{minipage}}\fi}}
\makeatother

%----------------------------------------------------------------------------------
%            content
%----------------------------------------------------------------------------------
\begin{document}
\nocite{*}
%-----       resume       ---------------------------------------------------------

\begin{tikzpicture}[remember picture, overlay]
  \node [anchor=north east, inner ysep=10mm, inner xsep=10mm]  at (current page.north east)
    {\includegraphics[width=75mm]{nvidia.png}};
\end{tikzpicture}

%\vspace{-2.5em}
\makecvtitle

%----------------------------------------------------------------------------------
%            university education
%----------------------------------------------------------------------------------
\section{Professional Experience}
	\cventry{since 2020}
					{Senior Deep Learning Compiler Engineer}
					{NVIDIA}
					{Würselen, Germany}{}
                    {
                        As part of NVIDIA's Deep Learning Compilation Group, I develop novel compilation techniques which will allow to generate high-performance programs for the rapidly changing landscape of specialized hardware accelerators.
					}
\section{University Education}
	\cventry{2016 -- 2020}
					{PhD degree in computer science}
					{University of Münster}
					{Münster, Germany}{}
                    {Supervisor: Prof. Sergei Gorlatch, Prof. Michel Steuwer (University of Edinburgh) \\
                        Thesis: \emph{High-Performance Domain-Specific Compilation without Domain-Specific Compilers.}
					}

	\cventry{2014 -- 2016}
                    {Master of Science in computer science (with distinction)}
					{}{}
					{}
					{University of Münster, Germany}

	\cventry{2011 -- 2014}
					{Bachelor of Science in computer science with a minor in mathematics}
					{}{}
					{}
					{University of Münster, Germany}
					{}
%----------------------------------------------------------------------------------
%            awards and grants
%----------------------------------------------------------------------------------
\section{Awards and Grants}
    \cventry{08/2020}{PhD thesis}{honored with the highest possible grade}{\strong{Summa cum laude}}{}{}
    \cventry{09/2019}{Heidelberg Laureate Forum 2019}{Invited Participation}{Invited to attend this highly prestigious forum with laureates of ACM A.M. Turing Award, Abel Prize, Fields Medal and Nevanlinna Prize. \textbf{Only 100 spaces} available for computer science}{}{}
    \cventry{12/2018}{NVIDIA Graduate Fellowship 2019}{50.000\$}{One of ten recipients world wide, selected as the only student from a European university}{}{}
    \cventry{02/2018}{Best Paper Award}{ACM International Symposium on Code Generation and Optimization (CGO) 2018}{Selected as the best paper out of 103 submissions}{}{}
    \cventry{since 2016}{Four Travel Grants}{HPC-Europa3, HiPEAC, 2x EuroLab-4-HPC}{total approx. 12.500\euro}{}{}

%----------------------------------------------------------------------------------
%            research visits
%----------------------------------------------------------------------------------
\section{Research Visits}
    \cventry{11/2019}
					{Visiting researcher (2 weeks)}
					{University of Glasgow}
					{Glasgow, UK}{}
					{During this visit, I prepared a publication about Elevate. We especially investigated how to implement TVM's scheduling language in terms of composable rewrite rules using Elevate.
					}


    \cventry{12/2018 \\-- 06/2019}
					{Deep Learning Compiler Engineer (7 months)}
					{NVIDIA}
					{W\"urselen, Germany}{}
					{
                    During this time, I extended the IR developed during my Internship in the US, focusing on code generation targeting NVIDIAs tensor cores in order to accelerate low-precision deep learning applications on GPUs.
                    }

    \cventry{07/2018 \\-- 09/2018}
					{Deep Learning Compiler Engineer Intern (3 months)}
					{NVIDIA}
					{Redmond, WA, USA}{}
					{
                    During this internship, I was working on an embedded DSL, IR and compiler for optimizing deep learning applications.
                    I focused especially on CUDA code generation for highly optimized matrix multiplication algorithms.
                    My work on this project simplified managing the complex compute and memory hierarchy on modern GPUs and enabled using program synthesis for design space exploration.
                    }

    \cventry{02/2018 \\-- 04/2018}
					{Visiting researcher (2 months)}
					{University of Glasgow}
					{Glasgow, UK}{}
					{Funded by HPC-Europa3\\
                     During this visit, I investigated the implementation of performance portable HPC applications with Lift including the automatic fusion of multiple compute kernels using rewrite rules.
                     I focused on geometric multigrid methods. The Irish Centre for High-End Computing (ICHEC) supported my visit by providing access to their GPU hardware.
					}

	\cventry{07/2017 \\-- 09/2017}
					{Visiting researcher (2 months)}
					{University of Edinburgh}
					{Edinburgh, UK}{}
					{Funded by HiPEAC\\
                     During this visit, I combined modern auto-tuning techniques with the current Lift code generator.
                     I also evaluated Lift's functional compilation approach compared to state-of-the-art polyhedral compilation.
                     A paper describing the results of this and our previous collaborations has won the \textit{best paper award} at the prestigious ACM International Symposium on Code Generation and Optimization (CGO)~\cite{cgo2018}
					}

	\cventry{02/2017 \\-- 03/2017}
					{Visiting researcher (2 months)}
					{University of Edinburgh}
					{Edinburgh, UK}{}
					{Funded by EuroLab-4-HPC\\
					 During this visit, I extended the Lift compiler, developed at the
					 University of Edinburgh, to enable automatic exploration of stencil-specific optimizations.
					}

	\cventry{04/2016 \\-- 05/2016}
					{Visiting researcher (2 months)}
					{University of Edinburgh}
					{Edinburgh, UK}{}
					{Funded by EuroLab-4-HPC\\
					 During this visit, I extended the Lift compiler to enable the generation of
					 high-performance stencil code for GPUs.
					}

	\cventry{09/2015}
					{Visiting researcher (3 weeks)}
					{HUST University}
					{Wuhan, China}{}
					{Funded by the EC’s 7th Framework Programme MONICA for accelerating the transfer
					 and deployment of research knowledge between European countries and China.
					 During this visit, I implemented an experimental setup for SDN-based multicast,
					 and prepared a research paper on this topic~\cite{humernbrum}
					}

%----------------------------------------------------------------------------------
%            publications
%----------------------------------------------------------------------------------
\printbibheading[title={Publications}]
\printbibliography[heading=none]
% Publications from a BibTeX file without multibib
%  for numerical labels: \renewcommand{\bibliographyitemlabel}{\@biblabel{\arabic{enumiv}}}% CONSIDER MERGING WITH PREAMBLE PART
%  to redefine the heading string ("Publications"): \renewcommand{\refname}{Articles}
%\bibliographystyle{plain}
%\bibliography{publications}

%----------------------------------------------------------------------------------
%            presentations
%----------------------------------------------------------------------------------
\section{Presentations}
    \cventry{08/2020}
					{\normalfont{Talk: \textit{Achieving High-Performance the Functional Way - A Functional Pearl on Expressing High-Performance Optimizations as Rewrite Strategies}}}{}{}{}
                    {ACM SIGPLAN International Conference on Functional Programming (ICFP 2020), Virtual Event}
    \cventry{03/2020}
					{\normalfont{Invited Talk: \textit{Elevate - A Language for Describing Optimization Strategies}}}{}{}{}
                    {Microsoft Research, Cambridge, UK}
    \cventry{01/2020}
					{\normalfont{Invited Talk: \textit{Elevate - A Language for Describing Optimization Strategies}}}{}{}{}
                    {NVIDIA, Redmond WA, USA}
    \cventry{04/2019}
					{\normalfont{Talk: \textit{High Performance Matrix Multiplication using Tensor Cores}}}{}{}{}
                    {NVIDIA, Redmond WA, USA}
    \cventry{03/2019}
					{\normalfont{Talk: \textit{Elevate - A Language for Expressing Optimization Strategies}}}{}{}{}
                    {Scottish Programming Language Seminar (SPLS), University of St. Andrews, UK}
    \cventry{12/2018}
					{\normalfont{Talk: \textit{High Performance Code Generation for Matrix Multiplication and Machine Learning}}}{}{}{}
                    {NVIDIA, Redmond WA, USA}
    \cventry{10/2018}
					{\normalfont{Talk: \textit{High Performance Geometric Multigrid Operations in Lift}}}{}{}{}
                    {HPC-Europa3 Transnational Access Meeting (TAM), Edinburgh, UK}
    \cventry{04/2018}
					{\normalfont{Talk: \textit{High Performance Stencil Code Generation with Lift}}}{}{}{}
                    {Workshop on Compilers for Parallel Computing (CPC), Dublin, Ireland}
    \cventry{04/2018}
					{\normalfont{Invited Talk: \textit{High Performance Stencil Code Generation with Lift}}}{}{}{}
                    {Dependable Systems Group, Heriot-Watt University Edinburgh, UK}
    \cventry{04/2018}
					{\normalfont{Tutorial: \textit{Lift: Performance Portable Parallel Code Generation via Rewrite Rules}}}{}{}{}
                    {International Symposium on Performance Analysis of Systems and Software (ISPASS), Belfast, UK}
    \cventry{03/2018}
					{\normalfont{Talk: \textit{High Performance Stencil Code Generation with Lift}}}{}{}{}
                    {Scottish Programming Language Seminar (SPLS), University of Glasgow, UK}
    \cventry{02/2018}
					{\normalfont{Talk: \textit{High Performance Stencil Code Generation with Lift}}}{}{}{}
                    {International Symposium on Code Generation and Optimization (CGO), Vienna, Austria}
	\cventry{02/2018}
					{\normalfont{Invited Talk: \textit{High Performance Stencil Code Generation with Lift}}}{}{}{}
					{Research Group on Compiler and Architecture Design, University of Edinburgh, UK}
    \cventry{03/2017}
					{\normalfont{Invited Talk: \textit{Performance Portable Stencil Code Generation with Lift}}}{}{}{}
					{Research Group on Compiler and Architecture Design, University of Edinburgh, UK}


%----------------------------------------------------------------------------------
%            research projects
%----------------------------------------------------------------------------------
\section{Research Projects}
    \cventry{}
                    {\normalfont{\small I have been actively contributing to the following research projects}}{}{}{}{}

    \cventry{since 07/2019}
					{Elevate}
					{\textit{A Language for Describing Program Transformations}}{}{}
                    {Ongoing research, \textit{https://github.com/elevate-lang}\\
					 I am leading the development of Elevate, a language enabling the specification of composable program transformations.
                     Using simple combination and traversal operators, Elevate allows to build domain-specific program optimizations as semantics preserving rewrite rules for arbitrary ASTs.
					}

    \cventry{since 09/2018}
					{Fireiron}
					{\textit{A Scheduling Language for High-Performance Linear Algebra on GPUs}}{}{}
                    {Ongoing research\\
					 I am leading the development of Fireiron, a scheduling language, IR and compiler for the generation of high-performance linear algebra kernels on GPUs.
                     Fireiron is a tool for experts, simplifying the development and experimentation with optimizations required for achieving highest performance on modern GPUs using specialized hardware units such as NVIDIA's Tensor Cores.
					}

	\cventry{2016 -- 2019}
					{Lift}
					{\textit{A Novel Approach to Achieving Performance Portability on Accelerators}}{}{}
					{\textit{www.lift-project.org}\\
					 I am one of the main contributors focusing on implementing stencil computations in Lift.
					 I extended the functional Lift IR and enabled the generation
					 of efficient OpenCL kernels for stencil-based applications.
					 The Lift project is a novel approach to generate high-performance OpenCL kernels
					 from high-level functional programs.
					}

	\cventry{04/2015}
					{PACXX}
					{\textit{Programming Accelerators with C++}}{}{}{
					 I developed an LLVM analysis pass for the PACXX compiler and ported HPC applications
					 to the PACXX programming model resulting in a publication~\cite{haidl}.
					 PACXX is a unified HPC programming model for programming accelerators
					 (GPUs etc.) using pure C++ by implementing a custom compiler
					 (based on the LLVM framework) and a runtime system.
					}
%	\cventry{10/2013 \\- 09/2014}
%					{OFERTIE EU Project}
%					{\textit{OpenFlow Experiment in Real-Time Internet Edutainment}}{}{}
%					{
%					 I configured the SDN
%					 testbed at the University of Münster, conducted several SDN-based experiments
%					 and extended the monitoring interface of the Real-Time Framework (RTF)
%					 The OFERTIE project aims to use SDN approaches to improve delivery of Real-Time Online
%					 Interactive Applications (ROIA).
%					}

%----------------------------------------------------------------------------------
%            reviewer
%----------------------------------------------------------------------------------
\section{Reviewer}

    \cvitem{2020}{CGO 2020 artifact evaluation chair}
    \cvitem{}{ASPLOS 2020 artifact evaluation committee}
    \cvitem{2019}{CGO 2019 artifact evaluation committee}
    \cvitem{2018}{CGO 2018 artifact evaluation committee}
    \cvitem{}{LCTES 2018 artifact evaluation committee}
    %\cvitem{2017}{Dr Hans Riegel Fachpreise commitee - (Undergrad Competition)}
	\cvitem{since 2016}
				 {I have been active as an external reviewer for the following conferences and journals:
                  \emph{Principles and Practice of Parallel Programming (PPoPP),
                    the International Parallel and Distributed Processing Symposium (IPDPS),
                    the International Journal of Parallel Programming (IJPP)},
					the \emph{Journal of Supercomputing},
					the journal \emph{Concurrency and Computation: Practice and Experience},
                    the Journal of Applied Geophysics (APPGEO),
					the Parallel Computing Technologies (PaCT),
					the Parallel Computing Conference (ParCo)
					%the UKRCON
					%and the PSI.
				 }

% %----------------------------------------------------------------------------------
% %            academic events
% %----------------------------------------------------------------------------------
% \section{Attended Academic Events}
% \cvitem{2020}{HiPEAC Conference, Bologna, Italy}
% \cvitem{2019}{Compiler and Programming Language Summit (Google), Munich, Germany}
% \cvitem{}{HLF - Heidelberg Laureate Forum, Heidelberg, Germany}
% \cvitem{}{ACACES Summer School (organized by HiPEAC) - \textit{Fifteenth International Summer School on Advanced Computer Architecture and Compilation for High-Performance and Embedded Systems}, Fiuggi, Italy}
% \cvitem{}{EuroLLVM - European LLVM Developers' Meeting, Brussels, Belgium}
% \cvitem{}{SPLS - \textit{Scottish Programming Languages Seminar}, St. Andrews, UK}
% \cvitem{2018}{Compiler and Programming Language Summit (Google), Munich, Germany}
% \cvitem{}{TAM - \textit{HPC-Europa3 Transnational Access Meeting}, Edinburgh, UK}
% \cvitem{}{CPC - \textit{20th Workshop on Compilers for Parallel Computing}, Dublin, Ireland}
% \cvitem{}{ISPASS - \textit{International Symposium on Performance Analysis of Systems and Software}, Belfast, UK}
% \cvitem{}{SPLS - \textit{Scottish Programming Languages Seminar}, Glasgow, UK}
% \cvitem{}{CGO - \textit{International Symposium on Code Generation and Optimization}, Vienna, Austria}
% \cvitem{2017}{Compiler and Programming Language Summit (Google), Munich, Germany}
% \cvitem{}{ACACES Summer School (organized by HiPEAC) - \textit{Thirteenth International Summer School on Advanced Computer Architecture and Compilation for High-Performance and Embedded Systems}, Fiuggi, Italy}
% \cvitem{}{PUMPS Summer School - \textit{Eighth edition of the Programming and Tuning Massively Parallel Systems summer school}, Barcelona, Spain}
% \cvitem{}{SPLS - \textit{Scottish Programming Languages Seminar}, St. Andrews, UK}
% \cvitem{2016}{HLPP conference - \textit{9th International Symposium on High-Level Parallel Programming and Applications}, Münster, Germany}
% \cvitem{}{UKMAC - \textit{UK Many-Core Developer Conference}, Edinburgh, UK}
% \cvitem{}{WadlerFest/LFCS30 - \textit{30th Aniversery of the Laboratory for Foundations of Computer Science}, Edinburgh, UK}
% \cvitem{2015}{PRACE Course - \textit{Advanced Parallel Programming with MPI and OpenMP}, Jülich, Germany}
% \cvitem{}{PRACE Course - \textit{Node-Level Performance Engineering}, Stuttgart, Germany}

%----------------------------------------------------------------------------------
%            memberships
%----------------------------------------------------------------------------------
%\section{Memberships}
%    \cvitem{}{ACM Student Member}
%    \cvitem{}{Berufungskomission W2 Praktische Informatik}

%----------------------------------------------------------------------------------
%            teaching
%----------------------------------------------------------------------------------
\section{Teaching}
	\cvitem{Winter '20}{Teaching assistant for the course: \textit{Parallel Systems}}
	\cvitem{Winter '20}{Supervised a student project: \textit{Developing a CUDA backend for Lift targeting TensorCores }}
	\cvitem{Winter '19}{Teaching assistant for the course: \textit{Operating systems}}
	\cvitem{Winter '19}{Teaching assistant for the course: \textit{Introduction to programming with Java and Haskell}}
	\cvitem{Summer '18}{Course design and Lecturer: \textit{Introduction to programming with C and C++}}
	\cvitem{Summer '18}{Teaching assistant for the course: \textit{Parallel Programming: Multi-Core and GPU}}
	\cvitem{Winter '17}{Teaching assistant for the course: \textit{Operating systems}}
	\cvitem{Winter '17}{Teaching assistant for the course: \textit{Introduction to programming with Java and Racket}}
	\cvitem{Summer '17}{Course design and Lecturer: \textit{Introduction to programming with C and C++}}
	\cvitem{Summer '17}{Supervised a student project: \textit{Automatic program optimization for modern many-core systems}}
	\cvitem{Winter '16}{Teaching assistant for the course: \textit{Operating systems}}
	\cvitem{Winter '15}{Student assistant for the course: \textit{Operating systems}}
	\cvitem{Summer '15}{Student assistant for the course: \textit{Computer architectures}}
	\cvitem{Winter '14}{Student assistant for the course: \textit{Operating systems}}

%----------------------------------------------------------------------------------
%            supervised students
%----------------------------------------------------------------------------------
\section{Supervised Undergraduate and Master Students}
    \cvitem{12/2018}{Johannes Lenfers (Master): \emph{Implementing Compiler Auto-Tuning Strategies for Design Space Exploration of Lift Programs}}
    \cvitem{02/2018}{Bastian Köpcke (Master): \emph{Efficient GPU Code Generation for FFT Computations in Lift}}
    \cvitem{01/2018}{Martin Lücke (Master): \emph{Efficient Implementation and Optimization of Geometric Multi-Grid Operations in the Lift Framework}}
    \cvitem{03/2018}{Clemens Hesse-Edenfeld (Undergraduate): \emph{Integrating Performance Models for Stencil Computations in Lift}}
    \cvitem{03/2018}{Alexander Dirk Holthaus (Master): \emph{Development of an Analytical Tool for Visualizing and Optimizing Memory Accesses in GPU Kernels}}
    \cvitem{03/2018}{Maurice Heine (Undergraduate): \emph{Implementation of a Visualization Tool for Lift Programs}}


%----------------------------------------------------------------------------------
%            technical skills
%----------------------------------------------------------------------------------
%\section{Technical Skills}
%	\cventry{Programming Languages}
%					{Scala, C/C++, Java}{}{}{}
%					{Experiences:
%					 Stencil support for Lift compiler (Scala),
%				 	 Multicast Module for the Floodlight SDN Controller (Java),
%					 Measurement library for OpenCL (C++),
%					 Implementation of the WiPo architecture (Java),
%					 Monitoring interface extension of RTF (C++)
%					}
%
%	\cventry{Parallel Programming}
%					{OpenCL, CUDA, OpenMP}{}{}{}
%					{Experiences:
%					 Performance portability evaluation of OpenCL Kernels on Intel Xeon (Phi) and NVIDIA Tesla.
%					 JIT compilation of a DSL using LLVM and CUDA Driver API
%					}
%
%	\cventry{Compiler Tools}
%					{LLVM}{}{}{}
%					{Experiences:
%					 Analysis Pass for the PACXX Compiler,
%					 Compiler frontend for self-defined DSL for data parallel applications based on algorithmic skeletons
%					}

\end{document}
